\subsection{KT-JOINT-RECHECK: the carried global-congruence reduction}
\label{sec:kt-global-residual}

The supplied checkpoint already proves K17-S17-JOINT,
K17-S17-GROUP-RING, and the joint operator test in
\texttt{evidence/work/hecke\_uniform/joint\_trace\_section.tex}.
The source audit below independently rechecks that carried
argument and makes its semisimplification consequence explicit;
it is not counted as a new arithmetic theorem. The argument
uses the lower-level model theorem underlying
\texttt{OneCharU\_s}, not merely its final quantified statement.
All archive inputs have the source-standing convention and the
pinned commit of KT-ONECHAR-CLOSED. No new Lean kernel build is
claimed. Write AV/ for the archive directory
\texttt{src/Development/AutomorphicArithmetic/AutomorphicVarieties/},
U3/ for AV/\texttt{UnitaryThreeGroups/}, and DTC/ for
AV/\texttt{ModularityArgument/ThresholdArguments/DenseThresholdCases/}.

Fix data $(\Omega,K,\ell,c,N,e,\psi,\iota)$ satisfying the
archive's \texttt{GD}, with unramified coefficient ring
$\mathcal O=c.O$; in particular $e(\mathcal O)=1$.
This includes the rational specialization $\mathcal O=\mathbf Z_\ell$.
Let $\mathscr A$ be the union of the two \emph{actual} admissible
classes in DTC/\texttt{NumberU3US.lean}:49--59:
\[
 \operatorname{Adm3U\_s}(\iota,e,N,\rho)
 \quad\hbox{or}\quad
 \operatorname{Adm2A}(\iota,e,N|\operatorname{disc}K|,\rho).
                                                               \tag{GR.1}
\]
Put $\Psi=\operatorname{Chr.Psi}(\ell,e,\operatorname{dirG}\psi)$
and write
\[
 \overline T=\overline\Psi|_{G_K}\oplus1
                    \oplus\overline\Psi^{-1}|_{G_K}.
                                                               \tag{GR.2}
\]
Residual representations are over $\overline{\mathbf F}_\ell$.
An integral model means matrices over the valuation ring of
$\overline{\mathbf Q}_\ell$ with inverse matrices also integral.

\begin{theorem}[K17-S17-JOINT, rechecked unramified specialization]
Under the nine closed archive inputs recorded in
KT-ONECHAR-CLOSED, if no $\rho\in\mathscr A$ has residual
semisimplification isomorphic to $\overline T$, then
\[
                   \operatorname{Vlen}(\Omega,c,e,\psi)=0.
                                                               \tag{GR.3}
\]
Equivalently, nonzero strict Selmer length forces one actual
admissible source representation with the global residual
semisimplification (GR.2). In fact the construction supplies
an integral model whose reduction is already the displayed
diagonal representation.
\end{theorem}

\begin{proof}
Use precisely the construction in
DTC/\texttt{FOneCharU.lean}:31--38: the closed inputs supply
$D,g_1,U,S,M$, where $S.R$ is the deformation ring and $M$
its graded model. They also supply finiteness of its
weight-zero subalgebra and the point hypothesis used in
\texttt{model\_upper}; the identical applications are written
at lines 55--57 of that file. These data are obtained before
using either the chosen Galois element or its trace hypothesis.

For a local point $f:S.R\to\overline{\mathbf Q}_\ell$ define
$\rho_f=D.\operatorname{rhoK}(\operatorname{push}f\,S.r)$.
Theorem \texttt{isPt\_push}, U3/\texttt{IsPtOf.lean}:124--166,
shows that the pushed representation has integral matrix
entries and, for every $g\in G_{\mathbf Q}$ and every pair of
indices $a,b$, satisfies
\[
 \left| (\operatorname{push}f\,S.r)(g)_{ab}
              -r_0(\eta,\Psi)(g)_{ab}\right|_\ell<1.
                                                               \tag{GR.4}
\]
Its entire proof of this clause is at lines 136--145:
$S.r.\operatorname{red}$ makes the difference lie in the
maximal ideal of $S.R$, and the defining local-point property
maps that ideal into elements of norm strictly less than one.
The integrality clause follows from \texttt{ipo\_le\_one},
lines 95--112. Applying it to $g^{-1}$ also gives integral
inverse matrices. These are the exact \texttt{int} and
\texttt{red} clauses of \texttt{Dat.IsPt},
U3/\texttt{RhoKApply.lean}:76--83.

On restricting to $G_K=\ker\eta$, the definition of
$D.\operatorname{rhoK}$ and its application lemma
(U3/\texttt{RhoKApply.lean}:70--74) identify these matrices
with those of $\rho_f$. Reduction of (GR.4) is therefore
exactly (GR.2). For every \texttt{TwoPt} point, the
\texttt{hadm} argument in DTC/\texttt{FOneCharU.lean}:47--54
places $\rho_f$ in (GR.1). Here \texttt{TwoPt} includes
\texttt{IsLocPt} by definition,
U3/\texttt{ConjStable.lean}:94--95.
Consequently the hypothesis of the theorem implies that
there are no \texttt{TwoPt} points at all.

Apply \texttt{model\_upper} with the identity Galois element.
Its sign condition is $\eta(1)=1$. Its remaining trace
condition says that the trace difference is nonzero at
every \texttt{TwoPt}; this is now vacuous. The theorem
(U3/\texttt{ModelUpper.lean}:104--112) yields finite model
length and a finite set $T$ satisfying
\[
 \operatorname{CoversPt}(T),\qquad
 \operatorname{nrm}(M.\operatorname{len})
       \leq2\sum_{f\in T}v_\ell(\operatorname{trDiff}(1,f)).
                                                               \tag{GR.5}
\]
By the first clause of \texttt{CoversPt}, every member of
$T$ is a \texttt{TwoPt}; see U3/\texttt{ConjStable.lean}:97--99.
Thus $T=\varnothing$ and the right side of (GR.5) is zero.
The normalized length is nonnegative; it is the finite
length divided by $e(\mathcal O)=1$, by
AV/\texttt{CharacterLevelBounds/ShaFamMono.lean}:71--83.
Hence $M.\operatorname{len}=0$. Finally
\texttt{vlen\_le}, U3/\texttt{VlenLe.lean}:52--67,
gives $\operatorname{Vlen}\leq M.\operatorname{len}=0$.
\end{proof}

\paragraph{Scope and use.}
The theorem supplies a necessary global congruence, not an
equivalence between a congruence and nonzero Selmer length.
The example KT-TWOFIELD-661 has genuine congruences while
the rational component vanishes. No ambient counting
formula alone rules out (GR.2). The reduction removes the
need to find a single Galois element separating every
ambient source, and it does not assume that the finite set
$X$ in \texttt{OneCharU\_s} is independent of that element.
It also keeps the \texttt{GD} and admissibility hypotheses:
it is not an assertion for arbitrary representations or
for an auxiliary field outside the established coverage.

\begin{proposition}[Horizontal form of the carried joint operator test]
Let $k$ be a field, $M$ a finite-dimensional $k$-space of
dimension $d$, and $A_1,\ldots,A_s$ commuting endomorphisms.
For specified scalars $t_1,\ldots,t_s\in k$, put
\[
 B=[\,A_1-t_1I\mid\cdots\mid A_s-t_sI\,].
                                                               \tag{GR.6}
\]
Then $B$ has row rank $d$ if and only if there is no
joint eigencharacter of the algebra $k[A_1,\ldots,A_s]$
on $M\otimes_k\overline k$ taking every $A_j$ to $t_j$.
Equivalently, the ideal generated by the $A_j-t_jI$ is
the unit ideal in that algebra.
\end{proposition}

\begin{proof}
Let $A=k[A_1,\ldots,A_s]\subseteq\operatorname{End}_k(M)$
and $I=(A_j-t_jI)_j$. The column space of $B$ is $IM$,
so full row rank is equivalent to $IM=M$.
The finite $A$-module $M$ is faithful by the definition of
$A$, and its support is therefore all of $\operatorname{Spec}A$.
For completeness, if $M_{\mathfrak m}=0$ at a maximal ideal,
finite generation supplies an element outside $\mathfrak m$
annihilating $M$, contradicting faithfulness. If $I$ were
proper, choose $\mathfrak m\supset I$. Nakayama applied to
the nonzero finite module $M_{\mathfrak m}$ gives
$IM_{\mathfrak m}\ne M_{\mathfrak m}$, a contradiction.
Thus $IM=M$ implies $I=A$; the converse is immediate.

If $I$ is proper, a maximal quotient of the finite
commutative algebra $A/I$ is a finite field extension of
$k$ and embeds in $\overline k$. This yields the required
character. Its localized module is nonzero by the same
faithfulness argument, so it occurs as a composition
factor of $M\otimes\overline k$ and, equivalently, as a
joint eigencharacter. Conversely such a character kills
$I$ and prevents $I=A$.
\end{proof}

In an established integral Hecke model covering all the
sources in (GR.1), apply this proposition to reductions
of commuting trace operators and the residual Eisenstein
traces. Full row rank for a finite collection of operators
rules out the global congruence (GR.2), even if every
individual determinant $\det(A_j-t_jI)$ vanishes.
Combining that certificate for every required source
branch with K17-S17-JOINT proves the desired vanishing.
The coverage and normalization remain explicit hypotheses;
a deficient block for finitely many primes does not by
itself establish a global congruence at all primes.

